\chapter{Introduction}

\section{Background}

Cloud-native applications are built from many small services packaged as containers
and orchestrated by \textbf{Kubernetes}. Kubernetes groups containers into
\emph{Pods}, schedules them on worker \emph{nodes} and isolates teams and applications
with \emph{namespaces}. Containers are not virtual machines: every container on a node
shares the same Linux kernel and is separated from its neighbours only by kernel
namespaces, control groups (cgroups) and capability restrictions. A process that
escapes this isolation, or that simply abuses the permissions it already has, can
affect every other workload on the node.

Most container security tools act \emph{before} a workload runs: images are scanned
for known vulnerabilities, manifests are checked against policies and admission
controllers reject risky configurations. These controls are necessary but they cannot
see what happens \emph{after} deployment. If an attacker exploits a vulnerable web
application and obtains a shell inside a pod, the actions that follow (downloading a
payload, opening a reverse shell, reading the service-account token, scanning the
cluster network or trying to break out of the container) are only visible at the level
of the operating system kernel.

The \textbf{extended Berkeley Packet Filter (eBPF)} makes this visibility practical.
eBPF allows small, verified programs to be loaded into the running Linux kernel and
attached to tracepoints, without kernel modules and without changing the application.
An eBPF program can observe every process execution, file open or network connection
on a node with very low overhead and stream the data to user space. Combining this
kernel-level telemetry with Kubernetes metadata and machine learning is the idea
behind this project.

\section{Motivation}

Runtime attacks on containers are common and well documented. Crypto-mining
campaigns such as Kinsing and TeamTNT compromise exposed containers and run miners
from hidden directories in \texttt{/tmp}; exploited web applications are used to open
reverse shells; stolen service-account tokens give attackers access to the Kubernetes
API; and privileged pods are abused to escape to the host with tools such as
\texttt{nsenter}. The MITRE ATT\&CK framework now contains a dedicated matrix for
containers that lists these techniques.

Existing open-source runtime tools (for example Falco, Tetragon and Tracee) are
powerful but are mainly rule-driven: they detect exactly what their rules describe and
require security expertise to tune. Pure machine-learning approaches, on the other
hand, are often evaluated offline on old system-call datasets and are rarely deployed
in a real cluster with Kubernetes context, latency measurements and an operator
interface. Small and medium-sized organisations in particular need a system that is
easy to deploy, explains \emph{why} something was flagged and is honest about its
detection rate and false positives.

This motivated us to build a complete, deployable runtime IDS that combines the
precision of rules with the generalisation of machine learning, enriches every alert
with Kubernetes context, and is evaluated not only on its own synthetic data but also
on real public datasets and live attack simulations.

\section{Problem Statement}

Kubernetes clusters lack an easy-to-deploy mechanism that can observe runtime
behaviour of containers at the kernel level, decide automatically whether that
behaviour is malicious, and present the result to an administrator with enough context
to act. In particular:
\begin{itemize}
    \item \textbf{No runtime visibility:} image scanning and admission control do not
          observe process, file and network activity after deployment.
    \item \textbf{Missing Kubernetes context:} kernel events carry PIDs and cgroup IDs,
          not Pod, container, namespace or workload names.
    \item \textbf{Rules alone do not generalise:} signature rules miss variants of an
          attack and need constant tuning.
    \item \textbf{ML alone is hard to trust:} a probability without a reason, severity or
          technique is difficult for an operator to act upon.
    \item \textbf{False positives:} legitimate platform behaviour (container runtimes,
          CoreDNS, health probes, CI jobs) often looks similar to attacks.
    \item \textbf{Overhead:} monitoring must not noticeably slow down the workloads it
          protects.
\end{itemize}

\section{Objectives of the Project}

The objective of this project is to design, implement and evaluate an eBPF-based
runtime intrusion detection system for Kubernetes. The specific objectives are:
\begin{enumerate}
    \item To capture security-relevant kernel events (process, file, network, privilege
          and namespace changes) from every container using \textbf{eBPF and libbpf
          CO-RE}.
    \item To map every event to its \textbf{Pod, container, namespace and workload}.
    \item To convert events into a fixed set of \textbf{behavioural features} that is shared
          by training and live inference.
    \item To build a labelled \textbf{runtime-telemetry dataset} including difficult edge
          cases, and to train a \textbf{Random Forest} classifier on it.
    \item To combine ML with \textbf{deterministic rules} into a 0--100 risk score with
          Low / Medium / High severity and a \textbf{MITRE ATT\&CK} technique.
    \item To store events and alerts in \textbf{MongoDB}, expose them through a secure
          \textbf{FastAPI} backend and visualise them in a \textbf{React} dashboard.
    \item To \textbf{evaluate} the system on accuracy, precision, recall, F1-score,
          false-positive rate, detection latency and resource overhead, on synthetic,
          real-world and live-cluster data.
\end{enumerate}

\section{Scope of the Project}

The system targets Linux nodes with kernel 5.8 or newer and BTF support, and was
developed and evaluated on a single-node \textbf{minikube} cluster. It \emph{detects}
and reports suspicious behaviour; it does not block or kill processes (prevention is
listed as future work). Detection covers nine attack categories: reverse shell,
dropper / payload, crypto-miner, privilege escalation, container escape, credential
access, discovery, defense evasion and lateral movement.

\section{Organisation of the Report}

\begin{itemize}
    \item \textbf{Chapter 2} surveys related work on system-call based intrusion
          detection, container security, eBPF monitoring and public datasets.
    \item \textbf{Chapter 3} analyses the existing and proposed systems and lists the
          functional, non-functional, hardware and software requirements.
    \item \textbf{Chapter 4} describes the datasets used for training and testing.
    \item \textbf{Chapter 5} explains the proposed method: sensing, enrichment, feature
          extraction, classification, rules and risk scoring.
    \item \textbf{Chapter 6} presents the system design with architecture, UML and data
          flow diagrams.
    \item \textbf{Chapter 7} describes the implementation, modules, deployment and the
          user interface.
    \item \textbf{Chapter 8} reports testing and the experimental results.
    \item \textbf{Chapter 9} concludes the report and outlines future work.
\end{itemize}
